2行1列のマス目が$n$個あると考える. このようなマス目について,
\begin{itemize}
    \item 2マスとも白であるものをA
    \item 2マスとも黒であるものをB
    \item 上が黒で下が白のものをC
    \item 上が黒で下が白のものをD
\end{itemize}
とよぶことにする. 条件を満たす塗り方をA,B,C,Dの4文字から(重複を許して)$n$個並べた文字列に翻訳すると,
\begin{center}
「AA」, 「BB」は部分列に含まれない
\end{center}
となるので, そのような文字列(以降\textbf{良い列}とよぶ)の個数$x_n$を求める.

良い列のうち, 右端がA,B,C,Dであるものをそれぞれ$a_n,b_n,c_n,d_n$とおく. 明らかに$x_n = a_n + b_n + c_n + d_n$. 長さ$n+1$の良い列$S_{n+1}$を右端とそれ以外の部分$S_n$に分割したとき, $S$は長さ$n$の良い列であって, 条件より
\begin{itemize}
\item $S_{n+1}$の右端がA なら $S_n$の右端はA以外
\item $S_{n+1}$の右端がB なら $S_n$の右端はB以外
\item $S_{n+1}$の右端がC,D なら $S_n$の右端はなんでもよい
\end{itemize}
これにより以下の関係式を得る:
\[
\begin{cases}
a_{n+1} &= b_n + c_n + d_n = x_n - a_n\\
b_{n+1} &= a_n + c_n + d_n = x_n - b_n\\
c_{n+1} &= x_n \\
d_{n+1} &= x_n
\end{cases}
\]
この4式の和を取り
\[
x_{n+1} = 4x_{n} - a_n - b_n
\]
一方で$n\geq 2$で$x_{n} = a_n + b_n + 2x_{n-1}$なので,
\[
x_{n+1} = 4x_n - (x_n - 2x_{n-1}) = 3x_n + 2x_{n-1}
\]
を得る($n\geq 2$). 以下, この漸化式を解く.

まず明らかに
\bealn{
x_1 &= 4 \\
x_2 &= 14
}
である. 特性方程式は$T^2 - 3T - 2 = 0$より
\[
\alpha = \dfrac{3 + \sqrt{17}}{2},\quad \beta = \dfrac{3-\sqrt{17}}{2}
\]
を用いて
\[
x_n = A\alpha^n + B\beta^n
\]
となるような定数$A,B$が存在する. 漸化式より$x_2 = 3x_1 + 2x_0$とすれば$x_0 = 1$と自然に定められるので, $n=0,1$を考え,
\[
\begin{cases}
A + B = x_0 = 1 \\
A\alpha + B\beta = x_1 =  4
\end{cases}
\]
より
\[
A = \dfrac{17 + 5\sqrt{17}}{34}, \quad B = \dfrac{17 - 5\sqrt{17}}{34}
\]
を得る.

以上より
{\scriptsize
    \bealn{
    &x_n = \\
    &\bolm{\parena{\dfrac{17 + 5\sqrt{17}}{34}}\parena{\dfrac{3+\sqrt{17}}{2}}^{n} + \parena{\dfrac{17 - 5\sqrt{17}}{34}}\parena{\dfrac{3 - \sqrt{17}}{2}}^n}
    }
}
